\subsection{Modelos de clasificación}

\subsubsection{Perceptrón multicapa}

Hemos usado el código de extracción de características para procesar los audios que tenemos en el dataset personalizado con 100 audios en español (50 reales y 50 generados a partir de esos reales). Una vez preprocesados los audios con sus características extraídas, empleamos una CNN (red neuronal convolucional) para sacar un embedding de 128 dimensiones del espectrograma del audio y un MLP (perceptrón multicapa) para sacar un embedding de 64 dimensiones de las características paralingüísticas del audio. Tras esto, las concatenamos teniendo un embedding de 192 dimensiones siendo las primeras 64 las características y las otras 128 las del espectrograma.

Después, creamos un perceptrón multicapa cuyo objetivo es clasificar el audio en si es generado o es real empleando para ello lógica difusa (\emph{fuzzy logic}). Para ello, pasamos la salida del perceptrón por un controlador fuzzy que irá aprendiendo también sus parámetros y adaptará la salida para que esté entre 0 y 1 siendo que cuanto más cercano al 0, más posible es que el audio sea real y no sea generado.
En la práctica, como el dataset es de 100 datos, no hay suficientes para entrenar un perceptrón de forma correcta ya que cuando hicimos las pruebas, todas las predicciones del perceptrón no se alejan del 0.5 que en este contexto de probabilidad continuo significa que el modelo está indeciso y no lo sabe clasificar como real o como generado.

\subsubsection{Fine-tuning Wav2Vec2}

Wav2Vec2 es un modelo de aprendizaje profundo para procesamiento de voz que permite obtener representaciones acústicas de alta calidad directamente desde la señal cruda de audio, sin necesidad de características manuales. Fue propuesto por Meta en 2020.

Este modelo recibe como entrada la forma de onda sin procesar del audio (audio crudo) y su salida debe ser decodificada utilizando Wav2Vec2CTCTokenizer.

El fine-tuning se hará con el dataset de 100 datos que tenemos para entrenar el Wav2Vec2 en la tarea de clasificación de audio. Al igual que con el modelo anterior, también lo pasaré por un controlador fuzzy para que nos diga cuan probable es que el audio sea real o IA.

% La parte de crear el dataset creo que sobra, eso ya se ha contado cómo se hace en el apartado del dataset
El proceso comienza creando un dataset de entrenamiento y otro de prueba en el que cada uno tendrá el mismo número de audios reales y generados y una etiqueta que indique si es humano o no. Una vez tenemos el dataset, hay que preprocesar los audios ya que en un inicio solo tiene las rutas de los audios por lo que ahora hay que cargar los audios de verdad. Para esta tarea usaremos el propio extractor de características que tiene el modelo (Wav2Vec2FeatureExtractor) fijando la frecuencia de muestreo a 16KHz y activando el padding que rellenará los audios cortos con ceros y trucará los audios largos para que todos tengan la misma longitud. Este preprocesado resultará en que ahora en el dataset, en lugar de haber rutas a los audios, hay tensores de Torch representándolos.

Ahora que ya tenemos el dataset de entrenamiento preprocesado, lo volvemos a dividir en entrenamiento y validación y entrenamos el modelo ya preentrenado para la nueva tarea de detección de voz generada por . Para evaluar el modelo se emplearon tres métricas, la precisión (verdaderos positivos / (verdaderos positivos + falsos positivos), la f1 (Media armónica entre precisión y recall) y el AUC (Área bajo la curva: Capacidad del modelo para distinguir entre clases positivas y negativas). El fine-tuning ha sido muy leve ya que con los 90 audios (80 de entrenamieno y 10 de validación) hay que congelar el aprendizaje del modelo para evitar un sobreentrenamiento.

Una vez hecho el fine-tuning, hacemos el mismo preprocesado para el dataset de test y ponemos el modelo en modo evaluación. Los resultados que de el modelo se pasan por un clasificador fuzzy que se encargará del postprocesado de la salida que la normalizará entre 0 y 1 siendo que los audios cuya score está entre 0.4 y 0.6 no los sabe clasificar, los que están por debajo de 0.4 son audios reales y los que están por encima de 0.6 son audios generados.

\subsubsection{Random Forest}

% El algoritmo Random Forest, propuesto por Breiman \cite{breiman2001random}, se seleccionó como uno de los clasificadores base debido a su robustez frente al sobreajuste, su capacidad para manejar características heterogéneas y su interpretabilidad mediante el análisis de importancia de variables.

% \subsubsection*{Justificación del modelo}

% Random Forest presenta varias ventajas para la tarea de detección de deepfakes de audio:

% \begin{itemize}
%     \item \textbf{Manejo de características heterogéneas}: Las características extraídas (formantes, MFCC, ZCR, HNR) operan en diferentes escalas y unidades; Random Forest no requiere normalización previa.
    
%     \item \textbf{Robustez ante el sobreajuste}: El mecanismo de bagging y la selección aleatoria de características reducen la varianza del modelo.
    
%     \item \textbf{Interpretabilidad}: Proporciona una medida intrínseca de la importancia de cada característica, permitiendo identificar qué parámetros acústicos resultan más discriminativos.
    
%     \item \textbf{Manejo de no linealidades}: Las firmas acústicas de las voces generadas por IA presentan patrones complejos que los árboles de decisión pueden capturar de forma natural.
% \end{itemize}

% \subsubsection*{Configuración experimental}

% El proceso experimental para Random Forest se estructuró en cuatro fases: (1) establecimiento de un modelo baseline, (2) validación cruzada para evaluar robustez, (3) optimización de hiperparámetros mediante búsqueda sistemática, y (4) evaluación comparativa con métricas exhaustivas.

% \paragraph{División de datos}
% El conjunto de datos se particionó en entrenamiento (80\%) y prueba (20\%) utilizando muestreo estratificado para mantener la proporción de clases en ambos conjuntos. Esta estratificación es crucial cuando existe desbalance entre voces reales y generadas.

% \begin{lstlisting}[language=Python, caption=División estratificada de datos]
% X_train, X_test, y_train, y_test = train_test_split(
%     X, y, 
%     test_size=0.2, 
%     random_state=42, 
%     stratify=y  # Mantiene proporcion de clases
% )
% \end{lstlisting}

% \paragraph{Modelo baseline}
% Se estableció una configuración inicial con hiperparámetros por defecto para obtener una referencia de rendimiento:

% \begin{minted}{python}
% rf_baseline = RandomForestClassifier(
%     n_estimators=100,      # Número de árboles
%     max_depth=10,           # Profundidad máxima
%     min_samples_split=5,    # Mínimo muestras para dividir nodo
%     min_samples_leaf=2,     # Mínimo muestras en hoja
%     max_features='sqrt',    # Características por división
%     random_state=42,
%     n_jobs=-1,              # Usar todos los cores
%     class_weight='balanced' # Ajuste para clases desbalanceadas
% )
% \end{minted}

% El parámetro \texttt{class_weight='balanced'} ajusta los pesos de las clases de forma inversamente proporcional a su frecuencia, compensando posibles desbalances en el conjunto de datos.

% \subsubsection*{Validación cruzada}

% Para evaluar la robustez del modelo y evitar que los resultados dependan de una única partición de los datos, se realizó validación cruzada estratificada con 5 folds:

% \begin{minted}{python}
% from sklearn.model_selection import StratifiedKFold, cross_val_score

% cv = StratifiedKFold(n_splits=5, shuffle=True, random_state=42)

% cv_scores_accuracy = cross_val_score(rf_baseline, X_train, y_train, 
%                                      cv=cv, scoring='accuracy')
% cv_scores_f1 = cross_val_score(rf_baseline, X_train, y_train, 
%                                cv=cv, scoring='f1')
% cv_scores_auc = cross_val_score(rf_baseline, X_train, y_train, 
%                                 cv=cv, scoring='roc_auc')
% \end{minted}

% La validación cruzada proporciona una estimación más realista del rendimiento esperado en datos no vistos, reportándose la media y la desviación estándar de las métricas.

% \subsubsection*{Optimización de hiperparámetros}

% Se realizó una búsqueda sistemática de hiperparámetros mediante \texttt{GridSearchCV} con validación cruzada de 5 folds, optimizando el F1-score. El espacio de búsqueda incluyó:

% \begin{itemize}
%     \item \texttt{n\_estimators}: Número de árboles en el bosque [50, 100, 200]
%     \item \texttt{max\_depth}: Profundidad máxima de cada árbol [5, 10, 15, None]
%     \item \texttt{min\_samples\_split}: Mínimo de muestras para dividir un nodo [2, 5, 10]
%     \item \texttt{min\_samples\_leaf}: Mínimo de muestras en una hoja [1, 2, 4]
%     \item \texttt{max\_features}: Número de características a considerar por división ['sqrt', 'log2']
% \end{itemize}

% \begin{minted}{python}
% param_grid = {
%     'n_estimators': [50, 100, 200],
%     'max_depth': [5, 10, 15, None],
%     'min_samples_split': [2, 5, 10],
%     'min_samples_leaf': [1, 2, 4],
%     'max_features': ['sqrt', 'log2']
% }

% grid_search = GridSearchCV(
%     RandomForestClassifier(random_state=42, class_weight='balanced'),
%     param_grid,
%     cv=cv,
%     scoring='f1',  # Optimización de F1-score
%     n_jobs=-1,
%     verbose=1
% )

% grid_search.fit(X_train, y_train)
% rf_optimized = grid_search.best_estimator_
% \end{minted}

% \subsubsection*{Métricas de evaluación}

% Para una evaluación exhaustiva de los modelos, se emplearon las siguientes métricas:

% \begin{itemize}
%     \item \textbf{Accuracy (Exactitud)}: Proporción de predicciones correctas sobre el total.
%     \[
%     \text{Accuracy} = \frac{TP + TN}{TP + TN + FP + FN}
%     \]
    
%     \item \textbf{Precision (Precisión)}: Proporción de predicciones positivas correctas.
%     \[
%     \text{Precision} = \frac{TP}{TP + FP}
%     \]
    
%     \item \textbf{Recall (Sensibilidad)}: Proporción de positivos reales correctamente identificados.
%     \[
%     \text{Recall} = \frac{TP}{TP + FN}
%     \]
    
%     \item \textbf{F1-Score}: Media armónica de precisión y recall.
%     \[
%     F1 = 2 \cdot \frac{\text{Precision} \cdot \text{Recall}}{\text{Precision} + \text{Recall}}
%     \]
    
%     \item \textbf{AUC-ROC} (Área bajo la curva ROC): Mide la capacidad del modelo para distinguir entre clases, independientemente del umbral de clasificación. Un valor de 0.5 indica clasificación aleatoria, mientras que 1.0 indica clasificación perfecta.
    
%     \item \textbf{MCC} (Coeficiente de correlación de Matthews): Métrica equilibrada que considera todos los elementos de la matriz de confusión, especialmente útil en problemas con clases desbalanceadas. Su rango es \([-1, 1]\), donde 1 indica predicción perfecta, 0 predicción aleatoria y -1 predicción inversa perfecta.
%     \[
%     \text{MCC} = \frac{TP \cdot TN - FP \cdot FN}{\sqrt{(TP+FP)(TP+FN)(TN+FP)(TN+FN)}}
%     \]
% \end{itemize}

% La función de evaluación implementada automatiza el cálculo de todas estas métricas y genera visualizaciones:

% \begin{minted}{python}
% def evaluacion_modelo(model_name, y_true, y_pred, y_pred_proba):
%     """
%     Evalúa un modelo con métricas completas y visualizaciones.
%     """
%     # Cálculo de métricas
%     accuracy = accuracy_score(y_true, y_pred)
%     auc = roc_auc_score(y_true, y_pred_proba)
%     f1 = f1_score(y_true, y_pred)
%     mcc = matthews_corrcoef(y_true, y_pred)
%     precision = precision_score(y_true, y_pred)
%     recall = recall_score(y_true, y_pred)
    
%     # Matriz de confusión
%     cm = confusion_matrix(y_true, y_pred)
    
%     # Visualizaciones (matriz de confusión, curva ROC)
%     # ...
    
%     return {
%         'accuracy': accuracy, 'precision': precision,
%         'recall': recall, 'f1': f1,
%         'auc': auc, 'mcc': mcc,
%         'confusion_matrix': cm
%     }
% \end{minted}

% \subsubsection*{Análisis de importancia de características}

% Una ventaja fundamental de Random Forest es su capacidad para proporcionar una medida de la importancia relativa de cada característica en la clasificación. Esta importancia se calcula como la reducción media de impureza (media de la disminución del índice de Gini) aportada por cada variable en todos los árboles del bosque.

% \begin{minted}{python}
% importances = rf_optimized.feature_importances_
% feature_names = X.columns if hasattr(X, 'columns') else [f'F{i}' for i in range(X.shape[1])]

% # Ordenar por importancia descendente
% indices = np.argsort(importances)[::-1]

% # Visualizar top 20
% plt.figure(figsize=(12, 8))
% top_n = min(20, len(feature_names))
% plt.barh(range(top_n), importances[indices[:top_n]][::-1])
% plt.yticks(range(top_n), [feature_names[i] for i in indices[:top_n]][::-1])
% plt.xlabel('Importancia relativa')
% plt.title('Top 20 características más importantes')
% plt.tight_layout()
% \end{minted}

% Este análisis permite identificar qué parámetros acústicos (determinados formantes, coeficientes MFCC, etc.) resultan más discriminativos para diferenciar voces reales de generadas por IA, aportando interpretabilidad al modelo y posibles hipótesis sobre las diferencias acústicas entre ambos tipos de voz.

% \subsubsection*{Resultados comparativos}

% La Tabla \ref{tab:rf_comparativa} muestra la comparación entre el modelo baseline y el modelo optimizado tras la búsqueda de hiperparámetros.

% \begin{table}[h]
% \centering
% \caption{Comparativa de rendimiento: Baseline vs Optimizado}
% \label{tab:rf_comparativa}
% \begin{tabular}{|l|c|c|c|}
% \hline
% \textbf{Métrica} & \textbf{Baseline} & \textbf{Optimizado} & \textbf{Mejora} \\
% \hline
% Accuracy    & 0.xxxx & 0.xxxx & +0.xxxx \\
% Precision   & 0.xxxx & 0.xxxx & +0.xxxx \\
% Recall      & 0.xxxx & 0.xxxx & +0.xxxx \\
% F1-Score    & 0.xxxx & 0.xxxx & +0.xxxx \\
% AUC-ROC     & 0.xxxx & 0.xxxx & +0.xxxx \\
% MCC         & 0.xxxx & 0.xxxx & +0.xxxx \\
% \hline
% \end{tabular}
% \end{table}

\subsubsection{LightGBM}

\subsubsection{Support Vector Machine (SVM)}

\subsubsection{AASIST}
Como hemos comentado antes, AASIST produce una predicción a partir del audio crudo por lo que el primer paso es definir las rutas a los audios. 

Después, se lee el contenido del \textbf{archivo .conf} del modelo. Para ello, hemos creado una función que busca el elemento pasado por parámetro y devuelve el bloque superior que contiene ese elemento junto con los elementos al mismo nivel. Esto nos permite construir una variable, que vamos a llamar args, con los siguientes campos:
\begin{itemize}
    \item \textbf{first\_conv}: configuración del \textbf{kernel} de la primera capa de convolución, es decir, cuántas muestras de audio se utilizan para cada filtro. En este caso, 128.
    \item \textbf{filts}: es un array de elementos siendo el primero el número de \textbf{filtros} en la primera capa de convolución. Los siguientes elementos son pares para cada bloque residual con el número de filtros de entrada y de salida.
    \item \textbf{gat\_dims}: tamaño del \textbf{embedding} que representa cada nodo. Tiene dos valores que hacen referencia al tamaño al construir los grafos de atención y al tamaño en las fases de ``HS-GAL'' + ``pooling'' respectivamente.
    \item \textbf{pool\_ratios}: se usa en la técnica de ``graph pooling'' y significa el \textbf{ratio de reducción de nodos} en cada capa. Los dos primeros valores se usan para reducir los grafos espectral y temporal respectivamente antes de combinarlos en el grafo heterogéneo. El siguiente valore se usa para reducir el grafo heterogéneo después de cada capa HS-GAL. Aunque existe un cuarto valor, en el modelo final no se usa.
    \item \textbf{temperatures}: controla lo suave o agresiva que es la \textbf{atención} en cada fase. Cuanto más bajo es este parámetro, mayor peso se concentra en pocos nodos vecinos.
\end{itemize}
Con args ya creado, la usamos para instanciar el modelo. Posteriormente, se usa el checkpoint para cargar los pesos preentrenados. 

El siguiente paso es definir las funciones de inferencia. Para esto, nos apoyamos primero en dos funciones auxiliares, una que convierte los audios a un \textbf{formato estándar}, un solo canal y 16 kHz, y otra que crea una \textbf{ventana} para recorrer audios largos. Esta última es importante ya que AASIST está diseñado para recibir segmentos de longitud de unos 4 segundos.
Con esto preparado, las tres funciones de inferencia las creamos para ver si conseguiamos mejores resultados alterando un un poco el proceso y se describen a continuación:
\begin{itemize}
    \item \textbf{average\_windows}: esta función devuelve la media de sumar las predicciones de cada ventana calculadas de forma independiente. 
    \item \textbf{score\_without\_windows}: esta función devuelve solamente la puntuación del segmento incial del audio en archivos largos o del audio en su totalidad en el caso contrario.
    \item \textbf{average\_logits}: esta función fusiona todas las ventana sumando sus logits antes de calcular una única probabilidad.
\end{itemize}
Finalmente, se aplica el modelo y se vuelcan los resultados en un csv con el nombre del audio, la probabilidad de que sea ``bonafide'' y su duración. Para esto, se recorre cada audio del dataset, se le aplica una de las funciones de inferencia y se escribe la información en el csv.